\section{Tangent cohomology and deformations}

We retain the standard filling data $(\ell _0,\ell _1,\ell _2)=(0,1,-1)$
and an arbitrary admissible $c$. All sheaves in this section are analytic,
and all cohomology groups have coefficients in $\C$. The integral twists
below mean the actual line bundles
\[
 M_n=\OO_X\bigl(n(S_1-S_2)\bigr),\qquad n\in\Z.
\]
We set $\mathfrak g=g_\gamma$. Thus, in the positively normalized
exponential coordinate of \ref{inv:C3.2},
$M_n=\exp(n\mathfrak g)$. The computations use the root base to retain the
nonreduced fibres. Passing to its coarse space is a separate, exact
invariant operation; it is never identified with taking the dual of an
invariant lattice.

\begin{proposition}[The root base and its explicit cohomology]\label{inv:D1.1}
Let $\rho:\mathcal B\to B$ be the analytic root curve with orders $3$ and
$4$ at $p_1,p_2$. Write $R_1,R_2$ for its tautological root line bundles.
Its Picard group is infinite cyclic, with generator
$R=R_1\otimes R_2^{-1}$, and
\[
 R_1=R^4,\quad R_2=R^3,\quad \rho^*\OO_B(1)=R^{12},
 \quad K_{\mathcal B}=R^{-7}.
\]
We abbreviate $R^n$ to $\OO_{\mathcal B}(n)$. Choose sections
$x\in H^0(R^3)$ and $y\in H^0(R^4)$ vanishing at $p_2$ and $p_1$,
respectively. For every integer $n$, put
\[
 r_1(n)\equiv n\pmod3,\quad 0\le r_1(n)<3,
 \qquad r_2(n)\equiv3n\pmod4,\quad0\le r_2(n)<4,
\]
\[
 d(n)=\frac{n-4r_1(n)-3r_2(n)}{12}.
\]
Then
\[
 \rho_*\OO_{\mathcal B}(n)=\OO_B(d(n)),\qquad
 H^i(\mathcal B,\OO(n))=H^i(B,\OO(d(n))).                 \tag{D.1}
\]
Equivalently, $h^0(\OO(n))$ is the number of pairs
$(a,b)\in\Z_{\ge0}^2$ with $3a+4b=n$, and
$h^1(\OO(n))=h^0(\OO(-n-7))$.
\end{proposition}
\Uses{the hypotheses, notation and definitions stated here; the explicit constructions and calculations in the proof}


The root-curve Picard group and graded cohomology are
\citet[Lemma~6.1]{DuEtAl2026v1}. We re-derive them and retain the explicit
coarse invariant-descent function $d(n)$ in the same calculation.

\begin{proof}
On the root charts the group acts on a disc coordinate by a primitive
$m_j$th root of unity. The stabilizer character of any line bundle is
removed uniquely by tensoring with $R_1^{-r_1}R_2^{-r_2}$, with the
exponents in the displayed ranges. A line bundle with trivial stabilizer
characters descends: an equivariant nonzero local frame can be averaged,
and its invariant frame gives the descent transition function. It is
therefore the pullback of a unique line bundle on $B$. Consequently
\[
 \operatorname{Pic}(\mathcal B)=
 \langle H,R_1,R_2\mid 3R_1=H,\;4R_2=H\rangle=\Z R.
\]
Here and below additive notation for this presentation means tensor
product. The determinant of the coordinate change $t=s^m$ gives
$K_{\mathcal B}=\rho^*K_B\otimes R_1^2R_2^3=R^{-24+8+9}$.

An invariant section of $R_j^{r_j}$ on its chart is $s^{r_j}$ times a
holomorphic function of $t=s^{m_j}$. It follows, by the local invariant
comparison on overlaps, that
$\rho_*(\rho^*\OO(k)R_1^{r_1}R_2^{r_2})=\OO(k)$.
In characteristic zero taking finite-group invariants is exact: the
averaging idempotent splits the invariant submodule. Thus there are no
higher direct images for $\rho$, and (D.1) follows. This proves the
cohomology comparison in the direction used here, not merely equality
of Euler characteristics. Finally every monomial of degree $n$ is
$x^{r_2}y^{r_1}$ times a homogeneous polynomial of degree $d(n)$ in
$x^4,y^3$. This gives the stated count of sections. The negative-degree
Cech complex on $x\ne0$, $y\ne0$ has cokernel represented by
$x^{-a}y^{-b}$ with $a,b\ge1$ and $3a+4b=-n$; replacing $(a,b)$ by
$(a-1,b-1)$ proves the $h^1$ formula without an additional duality
assumption. The two calculations agree exactly for every integer $n$;
the program \texttt{code/d\_exact.py} also compares them for
$-240\le n\le240$.
\end{proof}

\begin{proposition}[Vertical and logarithmic tangent bundles]\label{inv:D1.2}
There is a canonical representable proper flat holomorphic map
$\bar f:X\to\mathcal B$ with $f=\rho\bar f$. If
\[
 \mathcal T=T_X(-\log W),\qquad
 \mathcal V=\ker\bigl(\mathcal T\to
                   \bar f^*T_{\mathcal B}(-p_0)\bigr),
\]
then, for rank-two vector bundles $E,G$ on $\mathcal B$,
\[
 \mathcal V=\bar f^*E,\quad E=\OO\oplus\OO(-1),\qquad
 G=R^1\bar f_*\OO_X=\OO\oplus\OO(-1),                 \tag{D.2}
\]
\[
 R^0\bar f_*\OO_X=\OO,\qquad R^2\bar f_*\OO_X=\OO(-1),
\]
and the following is an exact sequence of locally free sheaves on $X$:
\[
 0\longrightarrow\bar f^*E\longrightarrow\mathcal T
 \longrightarrow\bar f^*\OO(-5)\longrightarrow0.       \tag{D.3}
\]
The trivial summands in $E,G$ are generated by the vertical field $\xi$
and the constant integral cocycle $\gamma$, respectively.
\end{proposition}
\Uses{the hypotheses, notation and definitions stated here; \cref{inv:D1.1}}


The root lift, vertical splitting, logarithmic quotient and
structure direct images are re-derived after reading
\citet[Lemmas~6.2 and 6.7; Propositions~6.6 and 6.10]{DuEtAl2026v1}.
The complete root and ordinary-sheaf arguments are retained below.

\begin{proof}
The divisors $S_j$ and their canonical sections give root data
$\OO_X(S_j)^{m_j}=f^*\OO_B(p_j)$, hence the indicated map to the root
curve. Over its order-$m_j$ chart the fibre product is the original
smooth torus family $J'_j$, since its free affine cyclic quotient is
exactly $f^{-1}(D_j)$. The comparison $J'_j\to X$ is finite etale;
in particular pullback is faithful and the stated local comparisons
descend. Away from the two root points the map is the original $f$.
Properness and flatness may be checked after this atlas: properness is
the properness of the original fillings, and flatness holds because a
nonconstant base parameter is a non-zero-divisor in every smooth
total-space local ring. Fibres have dimension two, including the cusp.

On a smooth torus chart a vertical field is a holomorphic linear
combination of the translation fields. Its coefficients are constant
on each compact fibre and therefore depend only on the base. The
vertical derivative of an affine transition is its linear part $R_g$;
the translating part, including the logarithmic regluing, has vertical
derivative the identity. At the cusp the two fields
$x_1\partial_{x_1},x_2\partial_{x_2}$ extend as logarithmic vector
fields in every toric chart. The deck maps multiply each $x_i$ by a
base-dependent unit and a power of the base parameter. Their vertical
logarithmic derivative fixes both coefficient directions. The
exponential overlap sends $\partial_{\zeta_i}$ to
$2\pi i x_i\partial_{x_i}$. These transitions depend only on the base
and define $E$; they prove the actual bundle identity
$\mathcal V=\bar f^*E$.

We check the kernel and surjectivity at every type of cusp point. At
$t=xy$, the kernel of $df$ is freely generated by
$x\partial_x-y\partial_y$ and $\partial_z$. At $t=xyz$ it is generated
by $x\partial_x-y\partial_y$ and $x\partial_x-z\partial_z$.
Indeed the equation $yz a+xz b+xy c=0$ in the regular local ring forces
$x\mid a,y\mid b,z\mid c$; after division the only relation is that
the three coefficients sum to zero. The logarithmic map sends
$\sum a_i z_i\partial_{z_i}$ to $(\sum a_i)t\partial_t$ and is
surjective. At a finite root point it is the ordinary relative tangent
sequence of the smooth map to $s$, not the singular differential of
$s^{m_j}$. These comparisons prove (D.3), since
$T_{\mathcal B}(-p_0)=\OO(7-12)$.

The invariant vector $e_2$ defines a subbundle $\OO\xi\subset E$.
It remains nonzero as a vector of the logarithmic vertical bundle at
the cusp, even where its image in the ordinary tangent fibre is zero.
The quotient line has automorphy factors $-1/\tau$ and $1/\tau$ under
the order-three and order-four generators. The modular discriminant
has weight twelve and no zero in the upper half-plane. Therefore its
reciprocal is a meromorphic section of the twelfth tensor power of
this quotient line. On the cusp chart
$\tau=(\log q)/(2\pi i)+h(q)$ with $h$ holomorphic, so the product
expansion of the discriminant is $q$ times a holomorphic unit.
There are no other zeros or poles, including on the finite root
charts. The quotient line thus has twelfth power
$\OO(-p_0)=\OO(-12)$, and the torsion-free Picard group in \ref{inv:D1.1} makes
it $\OO(-1)$. Its extension by $\OO$ splits because
$H^1(\mathcal B,\OO(1))=0$.

For completeness the structure-sheaf computation at the cusp uses the
normalization $P=\mathrm dP_6$, the three double curves and the two
triple points. Their exact normalization complex has global section
complex
\[
 \C\xrightarrow{0}\C^3
 \xrightarrow{(a,b,c)\mapsto(a-b+c,a-b+c)}\C^2.
\]
Exactness at the two local models is the elementary gluing of functions
on the coordinate branches; all higher cohomology of the rational
normalization, the rational double curves and the points vanishes.
The complex therefore gives $(h^0,h^1,h^2)=(1,2,1)$. Its product is
the Cech product of the two-triangle torus and gives a nondegenerate
alternating pairing $H^1\wedge H^1\to H^2$. This also follows directly
from the two opposite oriented triangles: the two coordinate cocycles
have determinant one on their oriented fundamental cycle. These are
the local comparison and product in Claude and Alpoge (2026,
Lemma 9.12), used here on the same normalization and with both triple
points retained.

On all other geometric fibres of $\bar f$ the cohomology is that of a
smooth two-torus. Constancy of these three dimensions and proper flat
analytic base change give locally free direct images on the atlas.
Their canonical comparison maps commute with the finite group actions
and with overlap restriction; hence they descend to $\mathcal B$.
The determinant of (D.3), using $W=\bar f^*p_0$, gives
\[
 \omega_{X/\mathcal B}=\bar f^*(\det E)^{-1}.
\]
Relative Serre duality, applied on that same atlas and descended by
its trace pairing, gives $R^2\bar f_*\OO_X=\det E=\OO(-1)$.
Cup product is an isomorphism $\det G\to R^2\bar f_*\OO_X$, since it
is so on every geometric fibre by the preceding calculation.

The class $\gamma$ is nowhere zero in $G$. On a smooth marked fibre,
$H^1(\OO)$ is the quotient of constant lattice cocycles by the two
period rows. If $\gamma$ were a combination of those rows, evaluation
on $\widehat w,\widehat\delta$ would make both coefficients zero,
contradicting $\gamma(\widehat\gamma)=1$. The same argument applies
on each finite root cover, and affine translations act trivially on
these classes. At the cusp $\gamma$ vanishes on the angular sublattice
and descends to a nonzero character of its rank-two deck lattice.
In the normalization complex it is the corresponding nonzero
one-cocycle on the two-triangle torus. Its cup product with a complementary
integral character is nonzero, so it remains nonzero there. Thus
$\OO\gamma\subset G$ is a subbundle; its quotient is $\OO(-1)$ by
the determinant just computed. Once again $H^1(\OO(1))=0$ splits the
extension. This proves every assertion in (D.2).
\end{proof}

\begin{proof}[Koszul and direct-image proof of Theorem~\ref{fam:st:A.5.2}]
Put $L=\OO_{\mathcal B}(-5)$. The first connecting map of (D.3) is
the holomorphic bundle map
$\kappa:L\to E\otimes G$. In the splittings fixed in \ref{inv:D1.2} the
coefficient in $\OO\xi\otimes\OO\gamma$ lies in $H^0(\OO(5))=0$.
The remaining coefficients lie in
$H^0(\OO(4))=\C y$, twice, and $H^0(\OO(3))=\C x$.
Therefore, as a map from $E^*$ to $G$ with its $L$ factor understood,
\[
 \kappa=\begin{pmatrix}0&a y\\ b y&c_0x\end{pmatrix}. \tag{D.5}
\]
We use $c_0$ for this coefficient to distinguish it from the deformation
parameter. The logarithmic period derivative at the cusp has leading
term the invertible matrix $B_0/(2\pi i)$. Hence
$\det\kappa=-ab y^2$ is nonzero there, so $ab\ne0$.
At the order-three point $\tau-\rho$ has order one on the root chart:
$j=1728t$ and the elliptic ramification order of $j$ at $\rho$ is
three. The marked torus period derivative is consequently nonzero
there: the two constant identity columns in the normalized period
matrix exclude an infinitesimal linear-frame change that could
otherwise conceal a nonzero derivative. In (D.5), $y=0$ and
$x\ne0$ at that point, so $c_0\ne0$.
These arguments establish the three nonzero coefficients; they do not
infer them from a global cohomology dimension. Rescaling the two
quotient frames and the source frame makes all three coefficients one.

The next connecting map is cup product with $\kappa$. In suitable
bases it is
\[
 L\otimes G=\OO(-5)\oplus\OO(-6)
 \xrightarrow{\ \delta_1\ }
 E\otimes\det G=\OO(-1)\oplus\OO(-2),\qquad
 \delta_1=\begin{pmatrix}-y&0\\-x&y\end{pmatrix}.       \tag{D.6}
\]
The sign follows by wedging first with $\gamma$ and then with the
second generator of $G$; reversing both sign conventions leaves the
cokernels unchanged. As its determinant is $-y^2$, it is injective as
a morphism of sheaves. The first connecting map is
$\delta_0=(0,y,y,x)^t$ and is a subbundle inclusion, since $x,y$ have
no common zero. Its cokernel is
\[
 \OO\oplus\OO(-1)\oplus\OO(2).                       \tag{D.7}
\]
Indeed one constant linear combination of the two $\OO(-1)$ summands
is unused, and the remaining quotient is given by the exact sequence
$0\to\OO(-5)\xrightarrow{(y,x)}\OO(-1)\oplus\OO(-2)
\xrightarrow{(x,-y)}\OO(2)\to0$.

The cokernel of (D.6) is supported at the order-three root point. There
$x$ is a unit, and elementary row and column operations identify it
with
\[
 Q=\OO(-1)\big|_{2p_1},                               \tag{D.8}
\]
where $2p_1$ means the length-two divisor on the root chart, not the
double coarse point. The generator is the first target component,
and its annihilator is $y^2$. The long direct-image sequence now gives
\[
 R^0\bar f_*\mathcal T=E,\quad
 R^1\bar f_*\mathcal T=\OO\oplus\OO(-1)\oplus\OO(2),
\]
\[
 0\longrightarrow Q\longrightarrow R^2\bar f_*\mathcal T
 \longrightarrow\OO(-6)\longrightarrow0.              \tag{D.9}
\]
This last extension splits, since $\operatorname{Ext}^1(\OO(-6),Q)
=H^1(Q(6))=0$; its splitting is not asserted canonical.

Let $\nu:P=\mathrm dP_6\to W\hookrightarrow X$ be the normalization
immersion. There is an exact sequence
\[
 0\longrightarrow\mathcal T\longrightarrow T_X
 \longrightarrow\nu_*K_P\longrightarrow0.             \tag{D.10}
\]
To verify the quotient, in a chart $W=(z_1\cdots z_r=0)$ the map takes
a vector field to its normal coefficient on each of the $r$ separate
branches. The coefficients lift independently; its kernel is precisely
the vector fields tangent to every branch. Thus the quotient is the
pushforward of the normal line of $\nu$, not $\OO_W(W)$.
The form $d\log x_1\wedge d\log x_2\wedge dq$ becomes a unit times
$dz_1\wedge dz_2\wedge dz_3$ in every unimodular chart and is invariant
under every cusp deck transformation. It trivializes $\nu^*K_X$.
Adjunction for the immersion gives its normal line $K_P$, proving
(D.10). Since $P$ is rational, its $K_P$ cohomology is $(0,0,1)$ by
Serre duality. Finite pushforward preserves that cohomology.

It remains to identify the extension at $p_0$ in degree two. We give
a fibre-sensitive calculation. The kernel of
\[
 \Omega_X^1|_W\longrightarrow\nu_*\Omega_P^1
\]
is $\nu_*\OO_P\otimes N_W^*$. At $xy=0$ its sections are
$a(y,z)y\,dx+b(x,z)x\,dy$; at $xyz=0$ they are
$a(y,z)yz\,dx+b(x,z)xz\,dy+c(x,y)xy\,dz$.
These identifications glue by the common conormal $dq$.
Since $N_W$ and $K_X|_W$ are trivial and $H^0(P,\Omega_P^1)=0$, the
Serre-dual space of $H^2(W,T_X|_W)$ is one-dimensional.
Top-degree base change identifies this group with the residue fibre
of $(R^2\bar f_*T_X)_{p_0}$. Its generic rank is one by (D.9).
A finite module over the discrete valuation ring $\C\{q\}$ of rank
one and residue dimension one is free of rank one. Thus the extension
by $\C_{p_0}$ increases the degree of its line part by $p_0$, without
adding torsion. We obtain the complete root direct images
\[
 R^0\bar f_*T_X=\OO\oplus\OO(-1),\quad
 R^1\bar f_*T_X=\OO\oplus\OO(-1)\oplus\OO(2),
\]
\[
 R^2\bar f_*T_X=\OO(6)\oplus Q.                       \tag{D.11}
\]
On the order-three chart, an invariant coefficient of the generator
of $Q$ must have exponent $2$ modulo $3$, whereas $Q$ only retains
exponents $0,1$. Therefore $\rho_*Q=0$. Applying the exact $\rho_*$
to (D.11) and using
$d(0)=0,d(-1)=d(2)=-1,d(6)=0$ proves (D.4).
Finally Leray over a curve has only columns zero and one. The usual
cohomology of these displayed line bundles gives the four asserted
dimensions. This is a direct-image computation, not an inference from
the tangent Euler characteristic.
\end{proof}

\begin{proof}[Actual-gluing Kodaira--Spencer calculation for Theorem~\ref{fam:st:A.5.2}]
All line bundles other than the trivial summands in
$R^q\bar f_*\mathcal V=E\otimes R^q\bar f_*\OO_X$ are acyclic:
they have weights $-1,-2$. Hence $H^1(X,\mathcal V)=
\C(\xi\otimes\gamma)$ and its other positive-degree groups vanish.
The three direct images of $\bar f^*\OO(-5)$ have weights
$-5$, $(-5,-6)$, and $-6$, all acyclic by \ref{inv:D1.1}. Consequently (D.3)
induces an isomorphism in every cohomology degree. Sequence (D.10),
whose quotient has no cohomology in degrees zero and one, then gives
the claimed isomorphism $H^1(\mathcal V)\to H^1(T_X)$.

We use the Cech convention in which the derivative of a transition
arrow is composed with the inverse of the central arrow. A lattice
transition is translation by $\Pi_c\lambda$, and its derivative in
$c$ is $\gamma(\lambda)\xi$. The linear deck matrices are independent
of $c$. On the finite charts the affine translations contribute
$\gamma(v_j)\xi/m_j$; their $m_j$th iterates give
$\gamma(v_j)\xi$, precisely the lattice cocycle. On the overlap with
the smooth part the derivative of the logarithmic translating section
is $\gamma(v_j)(\log s_j)\xi/(2\pi i)$. Its Cech coboundary accounts
for the difference of these representatives. On the cusp the deck
multiplier differentiates to $\gamma(\lambda)\xi$ in logarithmic
coordinates, and the exponential overlap identifies it with the same
cocycle. Thus the global Kodaira--Spencer class is
$[\xi\otimes\gamma]$, with a minus sign if the tangent-lift convention
$[\bar\partial h]$ is used instead. The sign is recorded; nonvanishing
does not depend on it.

Nonvanishing can also be checked on a regular torus. In the quotient
description of $H^1(\OO)$ used above, $[\gamma]\ne0$; the first
coefficient of (D.5) is zero, so $\xi\gamma$ cannot be a scalar
multiple of the base Kodaira--Spencer class. Its image in
$H^1(T_X|_F)$ therefore remains nonzero. Since $H^1(X,T_X)$ is
one-dimensional by \ref{inv:D1.3}, this proves the asserted isomorphism.
\end{proof}

\begin{proof}[Kuranishi ideal calculation for Theorem~\ref{fam:st:A.5.2}]
Let $\alpha$ be the generator in \ref{inv:D2.1}. The constructed holomorphic
family gives an actual convergent deformation tangent to $\alpha$.
Differentiating its integrability equation to second order gives
$\bar\partial\alpha_2=-[\alpha,\alpha]/2$, so the cohomology class
of $[\alpha,\alpha]$ is zero. Bilinearity proves that the bracket
$H^1(T_X)\times H^1(T_X)\to H^2(T_X)$ vanishes on every pair.
This argument alone would not exclude higher obstructions, so we
make the analytic argument as well. A Kuranishi germ is a closed
analytic subgerm of its tangent space $H^1(T_X)=\C$. Its classifying
map from the actual parameter disc has nonzero derivative by \ref{inv:D2.1}.
Composing with this embedding gives a holomorphic map of discs with
nonzero derivative; the inverse function theorem makes its image
contain a neighbourhood of zero in $\C$. Every equation of the
Kuranishi germ vanishes on that neighbourhood and is therefore zero.
The germ, including its analytic structure, is the smooth line.
The same inverse-function argument makes the classifying map an
isomorphism of germs. Pulling back the semiuniversal deformation
therefore realizes every sufficiently small deformation in the
$c$-family. This is semiuniversality; the positive-dimensional
automorphism group is not being discarded or interpreted as a
fine-moduli assertion.
\end{proof}




\section{Twisted tangent cohomology}
\subsection{All integral twists and the generic nonintegral twist}

\begin{proposition}[Integral Picard twists, including their torsion]\label{inv:D4.1}
For $M_n=\OO_X(n(S_1-S_2))$, define
$u(n)=\max\{d(n)+1,0\}$, $v(n)=\max\{-d(n)-1,0\}$ and
$\epsilon_n=0$ if $3\mid n$, $\epsilon_n=1$ otherwise. Then
\[
 f_*(T_X\otimes M_n)=\OO(d(n))\oplus\OO(d(n-1)),
\]
\[
 R^1f_*(T_X\otimes M_n)=
 \OO(d(n))\oplus\OO(d(n-1))\oplus\OO(d(n+2)),           \tag{D.12}
\]
\[
 R^2f_*(T_X\otimes M_n)=
 \OO(d(n+6))\oplus\C_{p_1}^{\epsilon_n}.
\]
In particular its cusp stalk is free of rank one. The four global
dimensions, for every integer $n$, are
\begin{align*}
 h^0&=u(n)+u(n-1),\\
 h^1&=v(n)+v(n-1)+u(n)+u(n-1)+u(n+2),\\
 h^2&=v(n)+v(n-1)+v(n+2)+u(n+6)+\epsilon_n,\\
 h^3&=v(n+6).                                          \tag{D.13}
\end{align*}
\end{proposition}
\Uses{the hypotheses, notation and definitions stated here; the explicit constructions and calculations in the proof}

\begin{proof}
The defining roots give $\bar f^*R=\OO_X(S_1-S_2)$, so the projection
formula applies to (D.11), with its actual coherent torsion summand.
Tensoring (D.8) changes its generator character from $-1$ to $n-1$.
There is exactly one invariant coefficient among exponents $0,1$ if
$n-1$ is $0$ or $1$ modulo three, and none if it is $2$. Multiplication
by the coarse parameter $t=s^3$ is zero on this module. Its invariant
direct image is therefore exactly the displayed skyscraper, of
length $\epsilon_n$, not a higher-order thickening. The two other
special places contribute no torsion. Formula (D.1) proves (D.12),
and the two-column Leray sequence proves (D.13). The code evaluates
these formulas for $-240\le n\le240$ using two independent integer
cohomology counts; their alternating sums are all one. That last
check is supplementary to the sheaf proof.
\end{proof}

\begin{proposition}[The cusp length and the generic twist]\label{inv:D4.2}
Use the coordinate $M_z=\exp(z\mathfrak g)$ of \ref{inv:C3.2}. For every
$z\notin\Z$ the cusp contribution is exactly
\[
 (R^0f_*(T_X\otimes M_z))_{p_0}=0,\quad
 (R^1f_*(T_X\otimes M_z))_{p_0}=0,\quad
 (R^2f_*(T_X\otimes M_z))_{p_0}=\C\{q\}/(q).           \tag{D.14}
\]
Outside the countable exceptional set explicitly defined in \ref{inv:D4.3},
these are the only nonzero direct-image stalks. Hence the generic
twist has $(h^0,h^1,h^2,h^3)=(0,0,1,0)$, and the cusp nilpotence exponent $r$ equals one there.
\end{proposition}
\Uses{the hypotheses, notation and definitions stated here; \cref{inv:C3.2}; \cref{inv:D4.3}; \cref{inv:D1.2}}

\begin{proof}
The map $H^1(W,\C)\to H^1(W,\OO_W)$ is an isomorphism by the
normalization complex in \ref{inv:D1.2}. The restriction of $M_z$ has trivial
pullback to $P$ and descent character
$(\exp(2\pi iz),1)$ on the two generators of the dual torus. The
normalization complex with these coefficients is the cellular
cochain complex of that rank-one local system. It may be computed
on the product cell structure as
\[
 \C\xrightarrow{(a-1,0)^t}\C^2
 \xrightarrow{(0,a-1)}\C,
 \qquad a=\exp(2\pi iz).
\]
For $z\notin\Z$ this is acyclic. All normalization components remain
rational with trivial pulled-back twist, so this calculates the
actual coherent cohomology $H^q(W,M_z|_W)$, not just a topological
analogy. Analytic semicontinuity and base change make
$R^q\bar f_*M_z$ zero on a neighbourhood of $p_0$ for all $q$.
The projection formula and (D.3) imply the same vanishing for
$\mathcal T\otimes M_z$. Tensoring (D.10) remains exact, and
$\nu^*M_z=\OO_P$ because $P$ is rational and the topological first
Chern class of $M_z$ is zero. Thus the degree-two direct image is
exactly the reduced skyscraper $H^2(P,K_P)\otimes\C_{p_0}$.
This proves (D.14), including its annihilator and not only its residue
dimension.

If the twist is nontrivial on every smooth fibre and on both finite
root-cover central tori, all cohomology of its fibre restriction
vanishes at those places. Indeed a nontrivial topologically trivial
line bundle on a complex torus has zero cohomology: in a unitary flat
presentation its Dolbeault Laplacian has no zero Fourier mode.
Tensoring the tangent extension by that line preserves acyclicity.
Base change on the root charts and exact invariants then give zero
direct images away from $p_0$. \ref{inv:D4.3} proves that the excluded twists
form a countable set. Leray now gives the asserted four dimensions.
\end{proof}

\begin{proposition}[The exceptional nonintegral twists cannot be omitted]\label{inv:D4.3}
On a regular marked fibre, $M_z$ is trivial if and only if there are
integers $n,p,q,r$ such that
\[
 q\tau(t)+r\mu(t)=p,\qquad
 z=n-6q\mu(t)-r\beta_c(t).                             \tag{D.15}
\]
For $z\notin\Z$ these resonances occur in a countable set of parameter
values $z$; every resonant $z$ is nonreal. On the finite root covers
the exceptional lattices are
\[
 \Lambda_1=\Z+(3\beta_c(z_1)+4-2\rho)\Z,\qquad
 \Lambda_2=\Z+(2\beta_c(z_2)+3-3i)\Z.                  \tag{D.16}
\]
Take the union of $\Z$, these two lattices, and all values from (D.15)
over regular fibres; this is the countable exceptional set used in
\ref{inv:D4.2}. The set is larger than $\Z$: the regular value $t=32/5$
supplies nonintegral twists with $h^1(X,T_X\otimes M_z)\ge2$.
\end{proposition}
\Uses{the hypotheses, notation and definitions stated here; \cref{inv:D4.2}}

\begin{proof}
The exponential sequence on a torus identifies trivial line bundles
with integral constant cocycles modulo the two holomorphic period
rows. Write
$z\gamma=(n,p,q,r)+a(6\mu,\tau,1,0)+b(\beta_c,\mu,0,1)$.
The third and fourth coordinates give $a=-q,b=-r$; the first two
give exactly (D.15). If $r=0$, the positivity of $\operatorname{Im}\tau$
forces $q=p=0$, leaving $z=n$. Otherwise the imaginary parts of
(D.15) give
\[
 \operatorname{Im}\mu=-\frac qr\operatorname{Im}\tau,
 \qquad\operatorname{Im}z=-rD\ne0.                    \tag{D.17}
\]
For fixed $(p,q,r)$ with $(q,r)\ne(0,0)$, the first equation of
(D.15) has discrete zeros, since $1,\tau,\mu$ are linearly
independent. A discrete subset of this second-countable Riemann
surface is countable. There are countably many integer tuples and
markings, so the union of the resulting $z$ is countable.

At $z_1$, insert $\tau=\rho$, $\mu=(2-\rho)/3$ in the first
equation. Its solutions have $(q,r,p)=(k,3k,2k)$.
At $z_2$ they have $(q,r,p)=(k,2k,k)$.
The second equation gives (D.16). These are discrete lattices because
their nonintegral generators have imaginary parts $3D$ and $2D$.

For the explicit regular resonance, use the section
$P=(0,\sqrt2t^2(t-1))$ of the stated elliptic surface. The
chord-and-tangent law gives
\[
 x(2P)=9t^2/8,\qquad
 y(2P)=t^2(32-5t)/(16\sqrt2).
\]
At $t=32/5$, the discriminant
$1728t^8(t-1)^3$ is nonzero, $P$ is not two-torsion, and $2P$ is a
nonzero two-torsion point. Thus $P$ has exact order four. The period
identification in Part A supplies a nonzero relation
$q\tau+4\mu=p$ with integers $p,q$. Equation (D.15), with $r=4$,
defines an actual twist for which the regular torus restriction is
trivial; (D.17) shows that it is nonintegral. The period derivative
there has nonzero determinant, so the tangent restriction has
$h^0(T_X|_F)=2$. The degree-zero base-change exact sequence gives
\[
 H^0(F,(T_X\otimes M_z)|_F)
   \simeq (R^1f_*(T_X\otimes M_z))_t[t],
\]
because $f_*(T_X\otimes M_z)=0$ for a twist nontrivial on the general
fibre. Its right side has dimension two. The degree-one direct image
is a finite-support sheaf, so its global sections have dimension at
least two, proving the bound on $h^1$. The exact elliptic calculation
is checked by both duplication and substitution in
\texttt{d\_exact.py}; its output records the nonzero rational
discriminant. The choice of integers $p,q$ depends on the marked
elliptic logarithm, and no unverified numerical marking is substituted.
\end{proof}

\begin{proposition}[Local complexes and every noncusp torsion length]\label{inv:D4.4}
Let $z\notin\Z$. At a regular resonance $t_*$ in (D.15), with
$A=\C\{t-t_*\}$, the two stalks are
\[
 (R^1f_*(T_X\otimes M_z))_{t_*}=(A/(t-t_*))^2,
 \qquad (R^2f_*(T_X\otimes M_z))_{t_*}=A/((t-t_*)^2).    \tag{D.18}
\]
At $p_1$, write $z=n+k(3\beta_1+4-2\rho)$, $k\ne0$, when this is
possible. Both coarse stalks are killed by $t$ and have length
$1+\epsilon_n$, with $\epsilon_n$ as in \ref{inv:D4.1}.
At $p_2$, write $z=n+k(2\beta_2+3-3i)$, $k\ne0$, when possible.
Both coarse stalks are killed by $t-1$ and have length one for
$n\equiv1,2\pmod4$ and length zero for $n\equiv0,3\pmod4$.
If the relevant lattice representation does not exist, both stalks
at that finite fibre are zero.
\end{proposition}
\Uses{the hypotheses, notation and definitions stated here; \cref{inv:D4.1}; \cref{inv:A2.3}}

\begin{proof}
We first give the comparison that turns the local calculation into
a computation of the sheaves. On a simply connected disc of smooth
torus fibres, including a finite root cover, retain the real marking
and the holomorphic period matrix. Translation by $\Pi(s)v$, for
fixed $v\in\R^4$, is holomorphic on the total family. After
subtracting the integer cocycle of a resonance, the line bundle has
constant lattice multipliers. Multiplying the lifted real translation
by its constant factor of automorphy makes the lattice act as the
identity and gives a genuine compact real torus action on the line
bundle. The same action acts on the tangent bundle by differentiation.
It commutes with the full Dolbeault differential.

Averaging over this compact torus projects the relative Dolbeault
complex onto its invariant part. Choose the lift so that the central
trivial line has trivial character. On every nonzero Fourier character,
the relative Dolbeault symbol is exterior multiplication by a
nonzero covector; contraction by a dual covector is a contracting
homotopy. After shrinking the disc, these nonzero covectors have a
uniform positive lower bound for bounded Fourier frequency and grow
linearly for unbounded frequency, because the real period matrices
and their inverses stay bounded. The homotopies therefore preserve
smooth Fourier series, with all derivatives, and are uniform on
compact subsets of the disc. The tangent extension is triangular
with respect to its vertical and horizontal parts; its nilpotent
off-diagonal term is treated by the finite homological perturbation
series, so every nonzero Fourier complex is still contractible.
Thus averaging is a quasi-isomorphism, not merely a fibrewise
dimension comparison.

The chain-level comparison can be written without an unspecified
holomorphic gauge. Write the normalized period matrix as $[Z(s),I]$
and use the real coordinates $\zeta=Z(s)x+y$, with
$x,y\in\R^2/\Z^2$. Set
\[
 G_j=(dx_j)^{0,1},\qquad e_i=\partial_{\zeta_i},\qquad
 h=\partial_s+\sum_i(Z'(s)x)_i e_i.
\]
These are invariant under the real translations. Since each $dx_j$
is closed on the total real product, $\bar\partial G_j=0$.
The forms $G_j$ frame the relative invariant $(0,1)$ forms because
$\operatorname{Im}Z$ is invertible, and
\[
 \bar\partial h=\sum_{i,j}Z'_{ij}(s)e_iG_j,
 \qquad (dy)^{0,1}=-Z(s)(dx)^{0,1}.
\]
For the smooth line frame defined by its flat character, the last
identity gives $\bar\partial\ell=2\pi i(a_1G_1+a_2G_2)\ell$,
with $a_1,a_2$ in (D.20); there is no additional base term.
All displayed coefficients are holomorphic in $s$. The full
invariant complex is thus the ordinary base Dolbeault resolution
of the following finite complex of holomorphic bundles:
\[
 C^j=(V\oplus L)\otimes\bigwedge^jG,\quad0\le j\le2,
 \qquad
 d(v\otimes\eta)=v\otimes a\wedge\eta,
\]
\[
 d(h\otimes\eta)=h\otimes a\wedge\eta+\kappa\wedge\eta. \tag{D.19}
\]
Here $V$ is the rank-two vertical bundle, $L$ the horizontal line,
$G=R^1\OO$ on the root chart, and $\kappa\in V\otimes G\otimes L^*$
is the actual tangent-extension cocycle. The matrix is obtained
by differentiating the period translations. In the raw period and
constant-cocycle frames it is
\[
 \kappa=\begin{pmatrix}6\mu'&\tau'\\\beta_c'&\mu'\end{pmatrix},
 \qquad
 a=(z-n+6q\mu+r\beta_c,-p+q\tau+r\mu).               \tag{D.20}
\]
The conventional factor $2\pi i$ in the Dolbeault representative of
the exponential class multiplies both components of $a$ by the
same nonzero constant; absorbing it in the frame has no effect on
any module below. The displayed formula also follows directly by
subtracting the holomorphic period rows from the constant lattice
cocycle. This identifies the chain map and its tangent extension;
no unspecified map between fibre cohomology groups is being used.
For the nonzero Fourier characters, use the filtration by the number
of base $(0,1)$ factors: the relative contractions make the associated
graded complexes acyclic, and the two-step filtration then proves
acyclicity of the full complexes. Thus the contraction need not
commute with the base Dolbeault operator. The finite cyclic group
normalizes the real translation group, so averaging and this
comparison are equivariant; exact invariants compute the actual
multiple-fibre stalks.

At a regular resonance, $a'=\kappa^t(q,r)^t$ and
$(q,r)\ne(0,0)$. The determinant of $\kappa$ is nonzero there by
Part A. Therefore $a$ has a simple zero. Choose its nonzero leading
direction as the second frame of $G$ and use invertibility of
$\kappa$ to choose the two vertical frames. Holomorphic invertible
changes reduce (D.19) to
\[
 a=(0,s),\qquad \kappa=\begin{pmatrix}0&1\\1&0\end{pmatrix}. \tag{D.21}
\]
The second cohomology is generated by $h\gamma\delta$ and is killed
by $s^2$; the relations eliminate the two vertical generators.
The first cohomology is generated by $e_1\delta,e_2\delta$, each
killed by $s$. These are equalities of modules over $\C\{s\}$
and prove (D.18).

There is an additional joint coefficient to check at the order-three
point. Put $t=s^3$, $\nu=\mu-(2-\tau)/3$ and
$b=\beta_c+4\mu+(2/3)\tau-b(0)$. The order-three transformation
laws give $\nu=\nu_2s^2+O(s^3)$ and $b=b_3s^3+O(s^4)$.
The elliptic calculation in \ref{inv:A2.3} gives, with elliptic scale
$f(s)=F s^2+O(s^3)$,
\[
 \nu_2^2=9F^2/512.
\]
For clarity we evaluate the other coefficients with the original
Weierstrass normalization. Since $x=f^2\wp$,
$y=f^3\wp'/2$, one has
$f^4g_2=12t^3(t-1)$ and $f^6g_3=-8t^4(t-1)^2$.
The identities $g_2=4\pi^4E_4/3$, $g_3=8\pi^6E_6/27$ and
$E_4'(\rho)=-2\pi iE_6(\rho)/3$ give
\[
 E_6(\rho)=-\frac{27}{\pi^6F^6},\qquad
 \tau(s)=\rho+\tau_1s+O(s^2),\quad \tau_1=\frac{i\pi F^2}{2}.
\]
The exact period determinant identity in Part A gives
\[
 K(s)=6\mu'^2-\tau'\beta_c'=K_2s^2+O(s^3),
 \quad K_2=\frac{3F^2}{4}.
\]
Equivalently $K=6\nu'^2-\tau'b'$, so
$K_2=24\nu_2^2-3\tau_1b_3$ and $b_3=7i/(32\pi)$.

For a nonintegral resonance the relation is
$(q,r,p)=\pm(k,3k,2k)$, $k\ne0$; the sign only changes the
following nonzero scalar. After a constant change of the two
cocycle frames, its Picard section is $3k(\nu,b)$, and the
transpose period derivative becomes
\[
 \begin{pmatrix}\tau'&\nu'\\6\nu'&b'\end{pmatrix}.
\]
Eliminating the unit $\tau'$ leaves the joint transverse component
$3k(b-6\nu\nu'/\tau')$. Its coefficient of $s^3$ is
\[
 -\frac{k(K_2+12\nu_2^2)}{\tau_1}
       =\frac{123ik}{64\pi}\ne0.                      \tag{D.22}
\]
Thus neither the order of the Picard section nor the determinant
order alone is used to select a Smith form. The joint coefficient
has been evaluated exactly.

Choose the second cocycle frame to be the Picard section divided by
$s^2$. Equivariance and (D.22) give a matrix with a unit in its
lower-right entry, first-order off-diagonal terms with nonzero
coefficients, and a second-order upper-left term. Subtracting a
holomorphic multiple of the second vertical row removes the last
term; the nonzero determinant coefficient and (D.22) keep both
off-diagonal coefficients units. Rescaling the equivariant frames
therefore reduces (D.19) to
\[
 a=(0,s^2),\qquad \kappa=\begin{pmatrix}0&s\\s&1\end{pmatrix}. \tag{D.23}
\]
Its cohomology is
\[
 H^1=\C\{s\}/(s^2)\oplus\C\{s\}/(s^3),\qquad
 H^2=\C\{s\}/(s^2)\oplus\C\{s\}/(s^3).             \tag{D.24}
\]
For example in degree one put
$u=e_1\delta$, $v=e_2\gamma+s h\delta$; the remaining cycle
$e_2\delta$ is eliminated by the boundary
$s u+e_2\delta+s v$. The two surviving relations are
$s^2u=0,s^3v=0$. In degree two eliminate $e_2\gamma\delta$;
the remaining generators $e_1\gamma\delta,h\gamma\delta$
have annihilators $s^2,s^3$ respectively.

We specify the character convention to avoid a sign ambiguity.
A frame has invariant exponent $w$ when its coefficient satisfies
$g(\zeta s)=\zeta^w g(s)$; its pullback-frame weight is the negative
of $w$. For the trivialized twist in (D.16), these exponents are
$n$ at $p_1$ and $3n$ at $p_2$. This agrees with $M_n$ in \ref{inv:D4.1}
and follows by evaluating its integer cocycle on $v_1,v_2$.
At $p_1$ the two generators in $H^1$ of (D.24) have invariant
exponents $n+2,n+2$, and the two generators in $H^2$ have exponents
$n+2,n$. The length-three summand has one invariant coefficient
for every $n$; the length-two summand has one precisely when
$n\not\equiv0\pmod3$. The coarse parameter is $s^3$, so all these
invariants are killed by it. This proves the asserted order-three
stalks with the full nilpotent layers included.

At $p_2$ the determinant is a unit. The Picard section has a simple
zero, so (D.21) applies on the order-four cover. The invariant
exponents of the two degree-one generators are $3n+1,3n+2$;
the cyclic degree-two module of length two has generator exponent
$3n+2$. In both degrees its invariant dimension is one exactly for
$n\equiv1,2\pmod4$. The coarse parameter $s^4$ kills the modules.
This proves the stated order-four calculation. The two independent
exact checks in \texttt{d\_koszul.py} are Smith normal forms and
ranks of truncated jet complexes; the latter are also restricted to
every stabilizer character. Their outputs agree with (D.21)--(D.24)
and with every invariant length just computed.
\end{proof}

\begin{proposition}[The full twisted cohomology and its jump loci]\label{inv:D4.5}
For $z\notin\Z$, let $\mathcal R_z$ be the set of regular base
points satisfying (D.15), counted as points, not as marked integer
tuples. It is finite. Put $N_z=|\mathcal R_z|$. Define
$\ell_1(z)=0$ off $\Lambda_1$, and
$\ell_1(z)=1+\epsilon_n$ if
$z=n+k(3\beta_1+4-2\rho)$, $k\ne0$.
Define $\ell_2(z)=1$ if
$z=n+k(2\beta_2+3-3i)$, $k\ne0$, with
$n\equiv1,2\pmod4$, and $\ell_2(z)=0$ otherwise.
Then, setting $J_z=2N_z+\ell_1(z)+\ell_2(z)$,
\[
 \boxed{(h^0,h^1,h^2,h^3)(T_X\otimes M_z)=(0,J_z,1+J_z,0).} \tag{D.25}
\]
More precisely, with $\mathfrak m_t$ the maximal ideal at $t$,
\[
 R^1f_*(T_X\otimes M_z)=
   \bigoplus_{t\in\mathcal R_z}\C_t^2
      \oplus\C_{p_1}^{\ell_1(z)}\oplus\C_{p_2}^{\ell_2(z)},
\]
\[
 R^2f_*(T_X\otimes M_z)=\C_{p_0}\oplus
   \bigoplus_{t\in\mathcal R_z}\OO_{B,t}/\mathfrak m_t^2
      \oplus\C_{p_1}^{\ell_1(z)}\oplus\C_{p_2}^{\ell_2(z)}. \tag{D.26}
\]
Together with \ref{inv:D4.1} these formulas cover every element of
$\operatorname{Pic}(X)=\C$. The actual cohomology jump loci are
locally discrete subsets of this complex line.
\end{proposition}
\Uses{the hypotheses, notation and definitions stated here; \cref{inv:D4.1}; \cref{inv:D4.2}; \cref{inv:D4.4}; \cref{inv:C3.8}}

\begin{proof}
The support assertions and all stalks were proved in \ref{inv:D4.2}--\ref{inv:D4.4}.
The direct image in degree zero is zero for a nonintegral twist:
it has generic rank zero and is torsion-free, since multiplication
by a nonzero base function is injective on the bundle on the smooth
integral total space. The support of either higher coherent direct
image is a proper closed analytic subset of the compact curve,
and hence is finite. This proves finiteness of $\mathcal R_z$ as
well as the absence of an accumulation at the special values.
If two integer tuples at a marked fibre represented the same
trivialized line, their difference would be an integral real class
in $H^{1,0}$, which is zero. Thus there is no multiplicity from
duplicate tuples. Formula (D.26) follows by the decomposition of a
finite-support sheaf into its local modules. Its global sections
have dimension equal to its length and its $H^1$ is zero, so Leray
gives (D.25).

To clarify the word generic, the global closed form $G_\gamma$ in
\ref{inv:C3.8} defines a holomorphic line bundle on $X\times\C$ by the
integrable operators $\bar\partial-2\pi i zG_\gamma$. Proper
coherent base change for the projection to $\C$ makes each
cohomology-dimension threshold an analytic subset. A real noninteger
$z$ has no resonance by (D.17) and (D.16), so the generic values are
attained. All possible exceptional values were shown to form a
countable set. Thus each actual threshold exceeding a generic
dimension is a proper analytic subset of a complex curve and is
locally discrete. Some indexed resonances, namely the two invisible
order-four congruences, produce no jump; (D.25) distinguishes them.

This is an explicit countable description of the jump loci in the
Picard coordinate, rather than a generic assertion. It can also be
indexed algebraically on the elliptic surface. If a regular fibre
has torsion section $P$ of exact order $r_t$, choose its marked
relation $q_t\tau+r_t\mu=p_t$. Its full resonance lattice is
\[
 \Z+\bigl(6q_t\mu(t)+r_t\beta_c(t)\bigr)\Z.
\]
Every regular resonance arises this way. The underlying values
$t$ are exactly the non-special roots of the division polynomials
$\psi_N(P)$, $N\ge2$. The nonintegral lattice points contribute two
to $J_z$ for each distinct such fibre. The finite-point lattices
have the congruence weights already stated. The script
\texttt{d\_division\_polynomials.py} evaluates the polynomials
exactly for $2\le N\le12$; for example the first nonconstant
regular factors are $5t-32$, $89t^2+4736t-4096$ and
$11t^2-6016t+8192$. These computations specify an executable
enumeration of the resonance loci without inventing decimal values
for the unspecified parameter $c$.

As an independent index check, $c_1(X)=c_2(X)=c_1(M)=0$ in integral
cohomology, so the degree-six Riemann--Roch integrand for the actual
twisted tangent bundle is $c_3/2$. Its integral is one. Both (D.13)
and (D.25) have alternating sum one. The index is used here to
verify the completed sheaf calculation, not to select its lengths.
\end{proof}

\paragraph{Source and re-derivation of the tangent theorem.}
The tangent cohomology $(1,1,1,0)$ and the smooth one-dimensional
Kuranishi germ are theorems of Du, Guo, Kusakabe, and Xie,
arXiv:2610.00250v1, Section~6 \citep{DuEtAl2026v1}.
The invariant run re-derived these results after reading that paper.
The present proof prints the root-base calculation, its local
comparisons, and the additional coarse direct images and all Picard
twists. Version~2 \citep{DuEtAl2026v2} does not contain the deformation
section; this does not alter the credit for the version~1 theorems.



\subsection{Tangent cohomology for arbitrary coprime multiplicities}\label{push:section:E.2}
\begin{theorem}[tangent cohomology and local miniversality]\label{push:src:E.2.1}
For every sphere member with multiplicities $(3k,4l)$ and the
admissible logarithmic data of \ref{push:section:E.1}, put $z=kl$ and $d=3k+4l$.
On $\mathfrak B_{k,l}=\mathbb P(3k,4l)$ the root lift
$\bar f:X\to\mathfrak B_{k,l}$ satisfies
\[
 \begin{split}
 \mathcal E&=\mathcal O\oplus\mathcal O(-z),\\
 R^q\bar f_*\mathcal O_X&=
 \mathcal O,\quad\mathcal O\oplus\mathcal O(-z),\quad\mathcal O(-z),\\
 K_X&=\bar f^*\mathcal O(z-d).
 \end{split}
\]
The tangent cohomology is $(h^0,h^1,h^2,h^3)(X,T_X)=(1,1,1,0)$.
With the groupoid derivative convention
$\kappa(\Phi)=(\partial_c\Phi)\Phi^{-1}$, the class of
$\partial/\partial c$ is $e\otimes[\gamma]$.  The class of
$\partial/\partial u$ is
$-(2\pi i u)^{-1}e\otimes[\gamma]$.  Both are nonzero.
The Kuranishi germ is a smooth reduced line, and the fixed-data
period family is locally miniversal.
\end{theorem}
\Uses{the hypotheses, notation and definitions stated here; \cref{push:section:E.1}; \cref{push:src:E.3.1}}

\begin{proof}
The root object is defined by the actual divisor identities
$f^*(0)=mS_1$ and $f^*(1)=nS_2$, with $m=3k$, $n=4l$.
It gives a holomorphic map to the analytic root stack.  Its base
changes to the two root discs are the actual smooth torus
families; the quotient actions are free.  At the cusp its charts
are $q=z_0z_1z_2$, so the source is smooth and the map is flat,
of relative dimension two, with Gorenstein fibres.  Properness is
checked on the root atlas, where it is properness of the compact
torus families and the original compact cusp filling.

Use weighted coordinates $x,y$ of weights $m,n$, with coarse
coordinate $t=y^m/(y^m-x^n)$.
The two root divisors have bundles $\mathcal O(n)$ and
$\mathcal O(m)$, and the cusp has bundle $\mathcal O(mn)$.
The root-forgetting morphism
\[
 r:\mathfrak B_{k,l}\longrightarrow\mathfrak B_{1,1},
 \qquad (x,y)\longmapsto(x^l,y^k),
\]
uses the group homomorphism $\lambda\mapsto\lambda^{kl}$ on
the weighted quotient.  It is a holomorphic morphism of analytic
stacks over $B$, with local maps $s\mapsto s^k$ and
$s\mapsto s^l$.  In particular $r^*\mathcal O(1)=\mathcal O(z)$.
For completeness, the analytic Picard group of this root curve has
the presentation
$\langle H,R_1,R_2\mid mR_1=H,\ nR_2=H\rangle$:
after subtracting the two stabilizer characters, a line bundle
descends to the coarse projective line.  Its only relations are
the two defining root identities.  The Smith form and
$\gcd(m,n)=1$ identify this group with $\mathbb Z$, with
$H=mn\mathcal O(1)$, $R_1=n\mathcal O(1)$, and
$R_2=m\mathcal O(1)$.  This proves the analytic torsion-free
Picard assertion used below.
Translations have identity vertical derivative and act trivially
on structure-sheaf cohomology.  Thus the vertical bundle computed
from the linear period factors is $r^*(\mathcal O\oplus
\mathcal O(-1))$.  The original determinant calculation can also
be repeated: its quotient line $L$ has
$L^{12}=\mathcal O(-p_0)=\mathcal O(-12z)$, so
$L=\mathcal O(-z)$ because the Picard group is torsion-free.
The fixed vector gives a subline $\mathcal O$; its extension
splits because $H^1(\mathcal O(z))=0$.

The cusp normalization resolution has terms
$\mathcal O_{\mathrm{dP}_6}$, three copies of
$\mathcal O_{\mathbb P^1}$, and two point sheaves.
Its global differential is $(a,b,c)\mapsto
(a-b+c,a-b+c)$, giving cohomology $(1,2,1)$; its alternating cup
pairing in degree one is nondegenerate.  The degree-one cocycle
$\gamma$ extends across the cusp because it vanishes on the toric
period lattice.  Its edge values $(0,-1,-1)$ have nonzero cup
product with $(1,0,-1)$.  On finite root charts its class is the
nonzero constant lattice cocycle $\gamma$, invariant under every
deck transformation.  These are explicit comparisons of the
same class, so it gives a nowhere-zero subline in
$G=R^1\bar f_*\mathcal O_X$.
The base-change and relative-duality argument \ref{push:src:E.3.1} applies on
these charts.  It gives $R^2\bar f_*\mathcal O_X=L$ and
$\det G=L$.  The extension of $L$ by the specified
$[\gamma]$-subline splits by the same $H^1(\mathcal O(z))=0$.
This proves all direct-image formulas.

The logarithmic tangent sequence is
\[
 0\longrightarrow\mathcal V=\bar f^*\mathcal E
 \longrightarrow\mathcal T=T_X(-\log W)
 \longrightarrow\mathcal H=\bar f^*\mathcal O(d-12z)
 \longrightarrow0.
\]
Its last map is the relative differential on smooth root charts
and sends $\sum a_i z_i\partial_{z_i}$ to
$(\sum a_i)q\partial_q$ at the cusp.  It is surjective there.
Since $K_{\mathfrak B}=\mathcal O(-d)$, determinants give
$K_X=\bar f^*\mathcal O(z-d)$.

We give the weighted cohomology calculation explicitly; simply
replacing $1$ by $kl$ in an acyclicity assertion would be incorrect.
The two quotient charts $x\ne0$, $y\ne0$ are finite quotients
of affine lines, and their overlap is a finite quotient of a
punctured line.  Finite averaging and Stein acyclicity compute
coherent cohomology by this two-chart complex.  Its homogeneous
Laurent monomials give
\[
 \begin{split}
 h^0\mathcal O(a)&=\#\{(i,j)\in\mathbb Z_{\ge0}^2:mi+nj=a\},\\
 h^1\mathcal O(a)&=h^0\mathcal O(-a-d).
 \end{split}
\]
The second equality identifies the remaining monomials
$x^{-i-1}y^{-j-1}$; it does not require a positivity theorem.
The bundles $\mathcal O(-z),\mathcal O(-2z)$ are acyclic.
Indeed a possible degree-one monomial would give
$m(i+1)+n(j+1)=az$, with $a=1$ or $2$.
Coprimality forces $i+1=lA$, $j+1=kB$, hence
$3A+4B=a$ with $A,B\ge1$, which is impossible.
Positive degree has no degree-one monomials, so
$H^1\mathcal O(z)=0$ as used above.

Put $b=d-12z$.  Both $b$ and $b-z$ are negative.
Degree-one monomials for $\mathcal O(b)$ would satisfy
$3A+4B=12$, with $A,B\ge1$, and none exist.
For $\mathcal O(b-z)$ the equation is $3A+4B=13$;
its only positive solution is $(A,B)=(3,1)$, producing
$x^{3l-2}y^{k-2}$ in the dual description.  It is allowed
exactly when $k>1$; write $e_k=1$ in that case and $e_k=0$
otherwise.  Projection formula and the two-column Leray sequence
now give
\[
 h^\bullet(\mathcal V)=(1,1,0,0),\qquad
 h^\bullet(\mathcal H)=(0,0,e_k,e_k),\qquad
 h^\bullet(\mathcal T)=(1,1,e_k,e_k).
\]
The program \texttt{e\_push\_exact.py} independently checks these
counts for all $1965$ coprime pairs with $k,l\le80$; the monomial
argument proves them for all pairs.

On the normalization $\nu:E=\mathrm{dP}_6\to W\subset X$,
the branchwise normal line is $K_E$: the invariant cusp form
$d\log x_1\wedge d\log x_2\wedge dq$ trivializes $K_X$
on the entire cusp neighbourhood.  Local normal coefficients
therefore give the exact sequence
\[
 0\longrightarrow\mathcal T\longrightarrow T_X
 \longrightarrow\nu_*K_E\longrightarrow0,
 \qquad h^\bullet(K_E)=(0,0,1).
\]
We must determine its connecting map when $k>1$.
The dual logarithmic sequence, tensored with $K_X$, has horizontal
subline $\bar f^*\mathcal O(13z-2d)$ and vertical quotient
$\bar f^*(\mathcal O(z-d)\oplus\mathcal O(2z-d))$.
The latter has no global sections, by the same equations
$3A+4B=1,2$ above.  Consequently
\[
 H^0(\Omega_X^1(\log W)\otimes K_X)
 =H^0(\mathfrak B,\mathcal O(13z-2d)).
\]
For $k>1$ this is generated by the monomial
$x^{3l-2}y^{k-2}$ multiplying the horizontal logarithmic
differential.  Both $x$ and $y$ are nonzero over $p_0$, so its
residue on $E$ is nonzero.  The ordinary logarithmic residue
sequence, tensored with $K_X$, identifies this residue map with
the Serre dual of
$H^2(E,K_E)\to H^3(X,\mathcal T)$.
One can check this identification directly in each chart
$q=z_0z_1z_2$: the residues of $dq/q$ are one on each branch,
and contraction with the invariant three-form is the normal
coefficient pairing.  Thus that connecting map has rank $e_k$.
The long exact sequence proves $h^2(T_X)=1$ and $h^3(T_X)=0$;
degrees zero and one are unchanged.

Finally $\partial_c\Pi=e\otimes\gamma$ gives the constant
groupoid cocycle $\lambda\mapsto\gamma(\lambda)e$ on every
regular torus.  On the enlarged finite roots the derivative of the
affine generator is $(r_j/m_j)e$, and the derivative of its
punctured connector is $(\log s_j)/(2\pi i)\,r_je$.
The identity $m_j(r_j/m_j)=\gamma(w_j)$ verifies the power
relation, while $\gamma A_j=\gamma$ verifies conjugation with
every lattice arrow.  Changing $\log s_j$ by $2\pi i N$
adds the derivative of the integral arrow $Nw_j$; the connector
formula is used only on the punctured overlap.  Its central
root-chart representative is the constant $(r_j/m_j)e$,
which is holomorphic.  At the cusp $\gamma$ kills the toric periods,
so the constant deck cocycle is exactly its extending cocycle
above.  Its image under the two isomorphisms
$H^1(\mathcal V)\simeq H^1(\mathcal T)\simeq H^1(T_X)$
is $e\otimes[\gamma]\ne0$.  The chain rule gives the asserted
$u$-class.  The Kuranishi local algebra is
$\mathbb C\{t\}/I$ since $h^1=1$.  The existing family induces
$t\mapsto g(c)$ with $g'(c_0)\ne0$.  The analytic inverse
function theorem makes this substitution injective; every
element of $I$ maps to zero, so $I=0$.  Its classifying map is
then an isomorphism of germs, proving local miniversality.
\end{proof}

\begin{proposition}[local separation and canonical dimension]\label{push:src:E.2.2}
The multiplicity pair $(3k,4l)$ is constant on every sufficiently
small deformation of one of the preceding members.  Distinct
pairs therefore determine disjoint open loci among these
isomorphism classes.  Moreover
\[
 \kappa(X)=\begin{cases}-\infty,&kl<3k+4l,\\1,&kl>3k+4l.
 \end{cases}
\]
Equality $kl=3k+4l$ never occurs under $\gcd(3k,4l)=1$.
\end{proposition}
\Uses{the hypotheses, notation and definitions stated here; \cref{push:src:E.1.2}; \cref{push:src:E.2.1}}

\begin{proof}
Local miniversality identifies every small deformation with the
fixed-data period family; \ref{push:src:E.1.2} makes its multiplicity pair an
intrinsic biholomorphic invariant.  The canonical formula of
\ref{push:src:E.2.1} and connected fibres identify its pluricanonical sections
with those of $\mathcal O(r(kl-3k-4l))$ on the root curve.
Negative weight has no sections.  Positive weight has linearly
growing section counts along multiples of $12kl$, giving
Kodaira dimension one.  Equality would imply $k\mid4l$;
coprimality forces $k=1$, and then $l=3+4l$ is impossible.
\end{proof}

\begin{proposition}[separation in irreducible analytic families]\label{push:src:E.2.3}
Suppose a proper holomorphic submersion over a connected
irreducible reduced complex base has two fibres of the type
in \ref{push:src:E.2.1}.  Then their multiplicity pairs agree, even if some
other fibres have not been constructed by this method.
For these members the asymptotic number
\[
 s(X)=\lim_{r\to\infty}
 \frac{h^0(X,K_X^r)-h^0(X,K_X^{-r})}{r}
 =\frac1{12}-\frac1{3k}-\frac1{4l}
\]
recovers the ordered multiplicity pair uniquely.
\end{proposition}
\Uses{the hypotheses, notation and definitions stated here; \cref{push:src:E.2.1}; \cite[Chapter III, Theorem 4.1]{BS76}}

\begin{proof}
For coprime $m,n$, the number of nonnegative solutions of
$mi+nj=N$ is $N/(mn)+O(1)$ as $N\to\infty$:
there is one permitted residue class of $j$ modulo $m$, and
its representatives lie between $0$ and $N/n$.
The canonical formula in \ref{push:src:E.2.1} proves the displayed limit,
with a negative value when the anticanonical counts grow.
Write $1/12-s(X)=a/b$ in lowest positive terms.
Since $\gcd(m+n,mn)=1$, we have $a=m+n$, $b=mn$.
Thus $m,n$ are the two roots of $T^2-aT+b$.
Their order is determined because only $m$ is divisible by
three and only $n$ is divisible by four.

For every fixed integer $r$, apply the finite-free-complex
theorem of \cite[Chapter III, Theorem 4.1]{BS76} to
the relative canonical power, which is an invertible flat
sheaf for a holomorphic submersion.  Its zeroth cohomology
dimension has analytic jumping loci.  On an irreducible
base any value constant on a nonempty open set is its
generic value: a proper analytic jumping locus has empty
interior.  Local miniversality in \ref{push:src:E.2.1} makes these dimensions
constant near either of the two given fibres, since all
nearby fibres have the same multiplicities and their
canonical formula is independent of the period parameter.
Consequently the two fibres have equal dimensions for
every positive and negative $r$.  Their limits $s(X)$
are equal, and the preceding arithmetic recovers the same
pair.  This proof uses relative cohomology only; it makes
no general invariance assertion for Kodaira dimension.
\end{proof}

\subsection{Analytic foundations and the version 1}\label{push:section:E.3}

\paragraph{Cited input.} \label{push:src:E.3.1}The source statement and its expanded derivation are retained in the proof compendium, Paper~I, label \nolinkurl{push:src:E.3.1}.\par

\Uses{the hypotheses, notation and definitions stated here; \cref{push:src:E.2.1}; \cref{fam:src:A.5.2}; \cite{DuEtAl2026v1}; \cref{app:rederive:push-src-E.3.1}; \cite[Section 6, Hauptsatz I, and Section 7, Satz 5]{Grauert1960D}; \cite{Grauert1963ErratumD}; \cite[Chapter III, Theorem 4.1]{BS76}}


This is analytic base change and duality, with the source hypotheses specified in the re-derivation. A complete re-derivation is given in Appendix~\ref{app:rederive:push-src-E.3.1}.


Grauert's original direct-image theorem
\cite[Section 6, Hauptsatz I, and Section 7, Satz 5]{Grauert1960D}
was also opened together with its correction \cite{Grauert1963ErratumD}.
The correction explicitly warns that the exceptional locus in
Section 7, Satz 2 need not be closed.  \ref{push:src:E.3.1} uses the
constant-cohomology theorem on every chart and the later
base-change formulation, so it does not make that erroneous
closedness inference.  For a generic argument over a smooth
curve one can instead use the finite free complex of
\cite[Chapter III, Theorem 4.1]{BS76}; its rank-jump minors
give the required local analytic exceptional sets.

\begin{proposition}[the five historical claims]\label{push:src:E.3.2}
For the standard twist vectors, each of the five named theorems
in the attached version 1 of \cite{DuEtAl2026v1} has a complete
argument relative to the explicitly constructed family and the
standard analytic theorems it cites.  We located no failed step
in those five arguments.  Their assertions are not present in
the attached version 2 \cite{DuEtAl2026v2}.  This textual
difference carries no inference about their mathematical truth
or about the reason for their omission.
\end{proposition}
\Uses{the hypotheses, notation and definitions stated here; \cite{DuEtAl2026v1}; \cite{DuEtAl2026v2}; \cref{push:src:E.3.1}; \cite[Theorem 7.4]{PY16push}; \cref{push:src:E.2.1}; \cref{fam:src:B.4.1}}

\begin{proof}
We give the dependency audit, rather than infer correctness from
the occurrence of the theorem labels.

For \texttt{p2:thm:tangent-cohomology}, the proof first constructs
the root lift as an actual Cartesian quotient, proves properness
and flatness on its atlas, and computes the cusp normalization
resolution with the multiplicative cup pairing.  The two lemmas
\texttt{p2:lem:analytic-grauert} and
\texttt{p2:lem:analytic-relative-duality} explicitly cite the
standard results checked in \ref{push:src:E.3.1} and perform descent; they do
not assume an unproved general stack duality principle.  Its
weighted cohomology and the normal sequence then give
$(1,1,1,0)$.  Its algebraic-to-analytic Picard comparison invokes
the applicable proper-stack GAGA theorem
\cite[Theorem 7.4]{PY16push}; the direct analytic Picard
argument in \ref{push:src:E.2.1} provides another justification.  Our general
calculation reduces to precisely its four acyclic weights
$-1,-2,-5,-6$ at $k=l=1$.  The extra logarithmic cohomology
for $k>1$ is outside that source's assertion.

For \texttt{p2:thm:kuranishi-aut}, the generating-arrow formulas
are $\gamma(\lambda)e$ on a lattice arrow, $0$ on a core deck
arrow, $e/3,-e/4$ on the two affine root generators, and
$\gamma(\lambda)e$ on a cusp arrow.  Its proof verifies the
relations, uses logarithms only on punctured connectors, and
uses the extending toric logarithmic vector on the cusp.
The Leray edge isomorphism carries this class to
$e\otimes[\gamma]\ne0$.  Completeness of the Kuranishi family
then gives a substitution $t\mapsto g(s)$ with $g'(0)\ne0$;
injectivity forces the entire defining ideal, including any
nilpotents, to vanish.  This proves a smooth reduced germ and
miniversality, without incorrectly claiming uniqueness or
universality.  The same theorem proves only
$\operatorname{Aut}^0=\mathbb C^*$: its flow has kernel
$\mathbb Z$, verified on the normalization of the cusp.  Full
automorphisms are treated separately below.

For \texttt{p2:thm:intrinsic-fibration}, adjunction and the
local finite characters give $K^{-2}=f^*\mathcal O_B(1)$.
Connected fibres and the projection formula identify all its
sections with the two base sections.  The complete pencil
therefore forces a biholomorphism to cover a base automorphism.
The three different multiplicity and monodromy types distinguish
the marked values, so this base automorphism is the identity.
This is the same argument as A.1 of the original return; it
needs no assertion that an arbitrary biholomorphism was
initially fibre preserving.

For \texttt{p2:thm:exact-period}, the simultaneous centralizer
is computed entrywise, and the period columns force integral
parameter difference.  Normalization after the genuine cyclic
base change tests the actual affine deck action, giving the
order-four parity obstruction.  Sufficiency is explicit:
choose $F^3=t$, $G^4=t-1$ and
$\ell=(2\pi i)^{-1}\log(F/G)$ on the universal cover, modulo
integral constants.  The source's translation is
\[
 b_{a,h}=-\frac h{12}\Pi_a\widehat\gamma+h\ell e.
\]
Its defects at $g_1,g_2,g_0$ are respectively
\[
 \tfrac h2(\widehat u-\widehat w-\widehat\delta),\qquad
 \tfrac h2(-\widehat w+\widehat\delta),\qquad0.
\]
For even $h$ these are integral.  Removing the known logarithmic
connector on each root chart leaves a holomorphic translation;
at the cusp $\ell=-s/12+$ a holomorphic function cancels the
period's logarithm.  Thus both the map and its inverse extend
over all three fillings.  Appendix calculations give the raw
composition defect $-hk\,e/12$ and the correction
$-h^2e/24$, which makes the action strict.  These identities
were re-executed exactly in
\texttt{code/e\_push\_source\_audit.py}.  The return's torsor
argument is another sufficiency proof, not a repair of a
missing proof in this source.

For \texttt{p2:thm:full-automorphisms}, a self-map's linear
part is first forced to be the identity.  The integral middle
translation cohomology is $\mathbb Z$, with representative
$c_N(g_1)=N(0,1,-2,0)^t$, $c_N(g_2)=0$.
The cusp defect is $N(0,-1,1,0)^t$, whose residual class in
$\Lambda/(\mathbb Z\widehat w+\mathbb Z\widehat\delta)$ is
$(0,-N)$.  A transverse disc through a smooth central point
lifts to the free cusp cover, forcing that residual monodromy
to vanish.  Hence $N=0$.  For the remaining zero class the
local quotient lift gives meromorphic character ratios; their
already single-valued additive logarithms have zero winding,
so the ratios are units and their logarithms extend.  Finite
averaging gives extension on both complete root discs.  The
result is a global section of the vertical bundle, hence a
constant multiple of $e$.  Its period kernel is $\mathbb Z$.
This proves $\operatorname{Aut}=\mathbb C^*$, including the
absence of disconnected components.  It corresponds to the
integral-cocycle and cusp-extension steps of \cref{fam:src:B.4.1}.

The audit program records attachment SHA-256 hashes and all five
label locations in version 1; every label and the corresponding
Kuranishi, Kodaira--Spencer, parity-classification and tangent
statements are absent from version 2.  We separately opened the
version-specific arXiv records and treat these as preprint
versions, not as a journal publication verified by its title.
\end{proof}

\begin{proposition}[limits of the overlap]\label{push:src:E.3.3}
The five audited version-1 results concern the standard
$(3,4)$ twist data.  That text does not classify general twist
vectors, compute the exact maximum $M$, construct a filling as
the family parameter tends to zero, compute a class in
$\Theta_7$, or settle the boundary questions in Parts A--D.
\end{proposition}
\Uses{the hypotheses, notation and definitions stated here; \cref{fam:src:C.1.1,fam:src:C.2.3}}

\begin{proof}
Its finite generators use only
$\widehat\gamma+2\widehat u-4\widehat w$ and
$-\widehat\gamma-3\widehat u+3\widehat w$ throughout the
deformation and classification arguments.  The cusp completion
constructed there fills the elliptic-base puncture for a fixed
period parameter; it is not a filling at zero of the parameter
$w$.  Its period domain is an admissible region, without the
maximization or the interval certificate of \cref{fam:src:C.1.1,fam:src:C.2.3}.  The concluding
orbit groupoid, quadratic correction, and the Oka results do not
supply the missing parameter degenerations or their topology.
Thus none of these omitted conclusions can be read into the
five theorems merely from the words ``family'' or ``cusp''.
\end{proof}

